\setcounter{section}{4}
\phantomsection
\section*{2. Literature Review}
\addcontentsline{toc}{section}{2. Literature Review}

In this section, we overview the key literature on birth order effects among historical scholars. We will look at four main categories. We will first summarise what is currently known about birth order and human capital formation in modern populations. Second, we discuss the mechanisms through which birth order may influence long-run outcomes. Thirdly, we will review what is known from historical and elite samples. Fourth, we will examine family institutions and inheritance. Lastly, we will discuss how this dissertation adds to the existing literature.

\phantomsection
\subsection*{2.1 Birth Order and Human Capital Formation}
\addcontentsline{toc}{subsection}{2.1 Birth Order and Human Capital Formation}

A large body of literature looks at birth-order effects in human capital accumulation. Theoretical frameworks such as the quantity-quality trade-off model developed by \cite{Becker1976} emphasized that parental resources (money, time, attention) are finite. Therefore, as family size increases, these investment resources are distributed among a larger number of children. This may reduce average investment per child and disadvantage later-born siblings. This framework is closely related to the resource dilution perspective developed by \cite{Blake1981}, which emphasizes how increases in family size may reduce investments per child.

Modern studies using administrative data and within-family comparisons find a robust ``firstborn advantage'' or ``firstborn premium'' in educational attainment and outcomes. Using Norwegian administrative data for cohorts born during the twentieth century, \cite{Black2005} show that later-born children complete approximately 0.2 to 0.3 fewer years of schooling. Because comparisons are made within families, these estimates are neither driven by differences in parental background nor socio-economic status across households.

Using U.S. family survey data from the late 20th century, \cite{Behrman1986} find that firstborn children obtain, on average, about 0.3 additional years of schooling. Using Canadian survey data from the late 20th century and early 21st century, \cite{Kantarevic2006} find that second-born children are 3 to 4 percentage points less likely to complete university education than firstborns, reinforcing the existence of birth-order differences in educational outcomes.

Similar results are found in other contexts. Using British household survey data from the early 2000s, \cite{Booth2009} show that birth order affects educational attainment. They show that firstborn children have higher levels of schooling and academic performance in comparison with their later-born siblings.

Taken together, these findings show that birth order is an important determinant of capital formation and a source of intra-family inequality. 

While this literature establishes a robust firstborn advantage in modern populations, it remains unclear whether these patterns extend to pre-industrial contexts characterized by higher fertility, different inheritance norms, and elite selection into education. We extend the birth-order literature to a historical population of scholars and examine whether the modern earlier-born siblings advantage is a historically persistent phenomenon.

\phantomsection
\subsection*{2.2 Mechanisms: Parental Investment and Sibling Dynamics}
\addcontentsline{toc}{subsection}{2.2 Mechanisms: Parental Investment and Sibling Dynamics}

There have been multiple proposals of mechanisms that might explain why birth order matters. One important explanation is differential parental investment. Using U.S. survey data from the early 2000s, \cite{Price2008} shows that ``a first-born child  receives 20-30 more minutes of quality time each day with his or her parent than a second-born child of the same age from a similar family''. These early differences accumulate over time and may have effects on skill formation.

Later research emphasizes strategic parental behavior. Using data from the late 20th century from the U.S., \citet{Hotz2015} show that parents respond to early indicators of skills and school performance by strengthening initial advantages rather than compensating weaker outcomes. Thus, relatively minor initial differences in cognition may increase with age, thereby adding to long-standing inequalities by birth-order.

Beyond parental investment mechanisms, there is a broader body of literature that emphasizes the importance of family background on individual outcomes. Studies using Scandinavian data on modern cohorts show that a large share of inequality can be attributed to factors shared within families, many of which are difficult to observe directly \citep{bjorklund2020intergenerational}. This perspective reinforces the relevance of within-family comparisons when studying birth-order effects.

Interaction with siblings also plays an important role. Later-born children grow up with less undivided parental attention and in environments where older siblings may act both as role models and as competitors. While psychological mechanisms are not the primary focus of economic models, these approaches are complemented by evidence such as \cite{Sulloway1996}, which argues that birth-order differences in non-cognitive traits, including risk-taking and openness to unconventional ideas, may persist over time.

In general, the literature identifies three principal channels through which birth order influences long-run outcomes: resource dilution; dynamic parental investment; and sibling competition.

Most existing studies focus on modern societies with formal schooling. This dissertation explores whether the same birth-order patterns can also be found among scholars from elite families in premodern times.

\phantomsection
\subsection*{2.3 Birth Order in Historical and Elite Contexts}
\addcontentsline{toc}{subsection}{2.3 Birth Order in Historical and Elite Contexts}

While the literature on historical birth-order effects is growing, it remains relatively limited and provides mixed evidence. The nature of fertility, mortality, educational institutions, and labor markets in pre-industrial societies was different from that of modern societies. These differences raise an important question: do the birth-order patterns documented in modern populations generalize to earlier periods?

Using large-scale census data for the United States between 1850 and 1940, \cite{cools2024birth} document stable firstborn advantages in occupational and social outcomes over time, despite major demographic and institutional changes. These findings indicate that birth-order effects can persist across periods of structural transformation, although their magnitude remains relatively modest.

Recent evidence from historical administrative data further supports the existence of birth-order effects in pre-industrial settings. Using data from the Netherlands during the 19th and early 20th century, \cite{nuevo2022birth} find a negative birth-order effect in occupational outcomes, with later-born individuals systematically attaining lower-ranked occupations than their older siblings. Importantly, they find that birth-order inequalities were already present before.

Recent contributions have further expanded the historical perspective by examining birth-order effects across the demographic transition and over very long time horizons. Using linked Dutch historical records from the 19th century, \cite{holthaus2025birth} find no systematic relationship between birth order and survival or longevity once individuals reach adulthood, with only small and context-specific effects in early-life outcomes. This suggests that, in pre-industrial environments characterized by high mortality and limited human capital accumulation, birth-order differences may have been muted or offset by other family and environmental factors.

Recent evidence using large-scale genealogical data suggests that birth-order effects may in fact be limited in historical contexts. \citet{ClarkCummins2024} analyze more than 400,000 individuals in England from 1680 to the present day and find a small impact of birth order on social outcomes in the general population. While a firstborn advantage is observed among a small group of elite families in earlier periods, birth order does not significantly affect outcomes for later-born children or across the broader population. These findings highlight that birth-order effects may be highly context-dependent and weaker in historical settings than in modern data. In this sense, focusing on scholars as a selected elite group may provide a setting where birth-order effects are more likely to emerge.

Taken together, these contributions suggest that birth-order effects in historical contexts are more heterogeneous than in modern data and may depend strongly on the institutional and demographic environment.

Likewise, \cite{DelacroixGoni2024} separate inherited human capital from nepotism using father-son links among European scholars from the Middle Ages to the Enlightenment. They demonstrate that while the intergenerational transmission of upper-tail human capital was strong, nepotism declined during the Scientific Revolution and Enlightenment. These results indicate that family background mattered in academia, but not all the mechanisms evolved through time as meritocratic pressures rose.

These studies document intergenerational persistence and the evolution of elite reproduction. But they do not systematically consider whether birth order within families influenced scholarly outcomes. It is an open empirical question whether the advantages that firstborns enjoy, as documented in modern data, were already present among historical elites. 

A small but growing literature using biographical data on European scholars between the Renaissance and the Industrial Revolution suggests that family background and early‑life conditions were strongly correlated with long‑run outcomes. Biographies of notable individuals are filled with systematic associations between family structure, longevity, and productivity; early family environments left marks that persisted even among highly selected cohorts \citep{DeLaCroix2013, DeLaCroixDocquier2024}.

Another recent research work has looked into the spatial and institutional factors behind elite creativity and intellectual performance. For instance, \cite{SerafinelliTabellini2022} demonstrate that the presence of creative elites in cities had persistent positive impacts for those Italian cities during the Renaissance and early modern period in terms of innovation and economic development. Conversely, the emergence of new large scale datasets on relevant historical individuals, like that constructed by \cite{LaouenanEtAl2022}, allows us to study the long-run rise and formation of European elites from the late Middle Ages to the Industrial Revolution.

This literature emphasises the importance of elite formation for long-run development, but does not consider intra-family dynamics like birth order. Whether elite trajectories were shaped by within-family hierarchy is directly relevant to this broader literature on the origins of creativity and upper-tail human capital.

Yet much of the existing work does not exploit within‑family identification or associate birth order with comprehensive long‑run achievement measures, which leaves open if birth-order patterns we observe today were already present among historical elites. Moreover, the mixed and often small magnitude of historical birth-order effects underscores the importance of examining specific elite populations and detailed productivity measures, where intra-family differences may be more likely to emerge.

Taking into account all this literature on historical elites and intergenerational transmission, this dissertation suggests that birth order may represent an additional intra-family dimension associated with elite formation, an aspect of family dynamics that has so far been largely neglected in scholarship.

This dissertation builds on the literature on historical elites and intergenerational transmission by examining the role of birth order within families. By comparing brothers and linking birth order to long-run measures of scholarly achievement, it provides new evidence on whether family hierarchy influenced academic success among historical elites.

\phantomsection
\subsection*{2.4 Family Institutions and Inheritance Systems}
\addcontentsline{toc}{subsection}{2.4 Family Institutions and Inheritance Systems}

One important dimension missing from the birth-order literature is that of family institutions. European regions between the sixteenth and nineteenth centuries exhibited different family systems, especially in their inheritance rules and household organization \cite{Todd1990}. In primogeniture systems, the eldest son inherited the assets; while in egalitarian systems, the assets were divided equally among the siblings.

These differences in family organization go beyond inheritance rules alone. As emphasized by \cite{BaudinDeRockGobbi2023}, family structures such as nuclear, stem, or complex families shape household decision-making processes, including investments in children and the allocation of resources within families. These institutional differences may therefore interact with birth order in determining individual outcomes.

Recent work indicates that customary institutions had long-term impacts on economic and demographic outcomes. For example, the Customary Atlas of Ancien Régime France covering French customary law before the Revolution from \cite{Gobbi2024} shows extensive local variation in customary law, including inheritance rules, marriage regulations, and parental authority rights. Such an exercise is useful to exploit spatial variation in family institutions for the identification of causal effects.

New studies using the Customary Atlas find that inheritance rules were closely tied to long‑run regional outcomes. The research demonstrates how customary inheritance systems influenced fertility behavior, regional development dynamics, income distribution, poverty and, household formation \citep{Gobbi2025}. Egalitarian inheritance regions were different from areas with primogeniture, suggesting a spatial variation in inheritance customs. Such family institutions have long-run implications for regional development, inequality, and social norms according to recent economic research \citep{Duranton2008, Gobbi2025}. 

In historical settings where the birth-order effects were likely influenced not only by biological or behavioral factors, but also by institutional context. Inheritance rules can maximize or minimize birth-order advantages, especially among elite families.

While recent research studies the long-run effects of inheritance systems on regional development, nothing is known about how these institutions determined inequality within families. This dissertation fills this gap by investigating whether regional inheritance regimes shaped within-family inequality among scholars, and in particular whether they reinforced or mitigated birth-order differences in long-run intellectual productivity.

\phantomsection
\subsection*{2.5 Contribution Relative to the Literature}
\addcontentsline{toc}{subsection}{2.5 Contribution Relative to the Literature}

Instead of being contributive along a single dimension, this dissertation connects across multiple strands of literature. It takes the birth-order literature to a historical sample of European scholars between the thirteenth and eighteenth centuries contributing to debates on elite formation and intergenerational transmission of privilege incorporating regional family institutions into intra-family inequality. This creates a single lens connecting birth order, institutions, and upper-tail human capital in the long-run.